\section{検出器}
\label{sec:detector}

本節は検出器を簡潔に述べる。
本研究の主張はアーキテクチャではなくデータ分布にあるため、
ここでは設計の要点と、
教師の整合性をどう保つかだけを記す。

\subsection{構成}
\label{sec:detector_arch}

入力画像を画像エンコーダへ通し、
パッチトークン列を得る。
本研究の実装では \BackboneName\cite{simeoni2025dinov3} を凍結して用いる。
タスク側の Transformer がパッチトークン間で自己注意を行い、
その出力を軽量な DPT 系デコーダ\cite{ranftl2021dpt}が画像解像度へ戻す。
密なヘッドは \KPName、領域、二値線、線の意味分類を予測し、
別のヘッドがカメラ姿勢と焦点距離を回帰する。

設計上重要なのは次の二点である。

\begin{enumerate}[leftmargin=1.6em]
  \item 密な予測と姿勢は\emph{並列}であり、
  検出した線を姿勢ソルバへ入力する直列構成ではない。
  したがって本研究は TVCalib や PnLCalib のソルバを内蔵しない。
  \item キーポイントのチャネルは物理点へ一対一に対応し、
  チャネル数の水増しや、
  画像座標からの順序推測を行わない。
\end{enumerate}

姿勢はカメラ中心の三値、
回転の連続六次元表現、
対数焦点距離の計十値として回帰する。
連続六次元表現\cite{zhou2019continuity}は
四元数の符号と正規化の曖昧性を避ける。
主点は補正後の対象内部パラメータから供給し、
ヘッドは予測しない。
焦点距離は等方かつ skew なしの契約であり、
レンズ歪みや主点移動は解かない。
この限界は\cref{sec:discussion} に記す。

\subsection{マルチタスク教師の整合}
\label{sec:task_consistency}

四種類の教師を同時に学習させると、
教師どうしの scope が食い違う危険がある。
本研究は次の整合条件をデータ側で強制する。

\begin{itemize}[leftmargin=1.6em]
  \item すべての密な教師と姿勢教師が\emph{同じ対象コート}を指す。
  多面の再構成では、
  対象コートを先に一つ決めてから全教師を生成する。
  \item キーポイントの物理番号が、
  密教師と姿勢教師で完全に一致する。
  食い違えばアダプタ構築時にエラーにする。
  \item 可視性と教師有効性のマスクが、
  全教師で同じ画像フレームに基づく。
  補正はすべての監督量を一度に更新する。
\end{itemize}

これらは性能を直接上げる仕掛けではなく、
教師の意味が一意に定まることを保証する条件である。
多面コートで教師の scope が混ざると、
損失は下がりながらコートの同一性が壊れる。
本研究はその状態を許さない。

\subsection{本稿が主張しないこと}
\label{sec:detector_scope}

複数のヘッドを持つこと自体は新規性ではない。
共有表現へ複数の幾何教師を与える研究は既に存在する\cite{nie2021robust}。
本研究が検出器について主張するのは、
再構成由来の教師が、
同じ幾何から生成されているがゆえに
教師間の整合条件を満たしやすいという点と、
その条件を実装で強制しているという点である。
アーキテクチャの優劣や勾配の診断は本稿の範囲外であり、
必要な比較は\cref{sec:experiments} の計画に含める。
